\FifthArticle{15}
Über die Anwendung dieser Vorschriften ist noch Folgendes hinzuzufügen.

Zur Bestimmung jeder Wurzel sind zwei Operationen auszuführen: Zuerst ist aus $\lambda$ der zugehörige Wert von $a$ (und damit zugleich der von $b$ oder $c$) abzuleiten; dann ist daraus der Wert von $x$ zu berechnen. Für jede dieser beiden Operationen kann man zwischen drei Formeln wählen; ich bevorzuge in den meisten Fällen die zuerst angegebenen. Bei allen diesen Rechnungen hat man es gar nicht mit den Größen $\lambda$, $a$, $b$, $c$ selbst zu tun, sondern nur mit ihren Logarithmen. Die erste Operation ist indirekt und beruht daher in der Regel auf mehreren schrittweise fortschreitenden Näherungen. Dabei wird es bequem sein, anfangs Tafeln mit einer geringeren Ziffernzahl zu verwenden. \textsc{Matthiessen}s Tafel hat bekanntlich sieben Dezimalstellen, meine fünf; \textsc{Encke} und \textsc{Ursin} haben sie mit vier Ziffern abdrucken lassen. Wenn man zu Beginn der Arbeit noch gar keine Kenntnis einer ersten groben Näherung mitbringt, wird man es vielleicht vorteilhaft finden, einen noch kürzeren Auszug aus den Tafeln mit nur drei Ziffern auf einem eigenen kleinen Blatt vor sich zu haben, etwa so:
\[
\begin{array}[t]{c|c|}
A&B\\\hline
0&0{,}301\\
0{,}1&0{,}254\\
0{,}2&0{,}212\\
0{,}3&0{,}176\\
0{,}4&0{,}146\\
0{,}5&0{,}119\\
0{,}6&0{,}097\\
0{,}7&0{,}079\\
0{,}8&0{,}064\\
0{,}9&0{,}051\\
1{,}0&0{,}041
\end{array}
\begin{array}[t]{c|c|}
A&B\\\hline
1{,}0&0{,}041\\
1{,}1&0{,}033\\
1{,}2&0{,}027\\
1{,}3&0{,}021\\
1{,}4&0{,}017\\
1{,}5&0{,}014\\
1{,}6&0{,}011\\
1{,}7&0{,}009\\
1{,}8&0{,}007\\
1{,}9&0{,}005\\
2{,}0&0{,}004
\end{array}
\begin{array}[t]{c|c}
A&B\\\hline
2{,}0&0{,}004\\
2{,}1&0{,}003\\
2{,}2&0{,}003\\
2{,}3&0{,}002\\
2{,}4&0{,}002\\
2{,}5&0{,}001\\
2{,}9&0{,}001\\
3{,}0&0{,}000
\end{array}
\]
\SourceLeaf{100}

